\subsection{Quantitative convergence of the currents}

The bound for the constants $c_\lambda$ in \cref{prop:semipositive-dual-exponential-limit} also bounds the trace error linearly in $\lambda$. Away from the degeneracy set, the deficit in the arithmetic--geometric mean inequality converts this bound into a rate of convergence for all coefficients of the currents. We keep the notation $u_0$, $v_\lambda$, $c_\lambda$ of that proposition. The total variation of a current of order zero is measured with respect to a fixed smooth Hermitian metric; on a relatively compact chart it is equivalent to the sum of the total variations of its coefficients.

\begin{proposition}\label{prop:semipositive-dual-current-rate}
Under the hypotheses and notation of \cref{prop:semipositive-dual-exponential-limit}, put
\[
 A_\lambda=\omega+\ddc v_\lambda,\qquad
 A_0=\omega+\ddc u_0.
\]
There is a constant $C$ such that
\[
 \|A_\lambda\wedge T-A_0\wedge T\|_{\mathrm{TV}}
       \leq C\lambda\qquad(\lambda>0).
\]
For each compact set $K\Subset X\setminus Z$, there is a constant $C_K$ such that
\[
 \|A_\lambda-A_0\|_{\mathrm{TV}(K)}
       \leq C_K\sqrt\lambda\qquad(0<\lambda\leq1).
\]
\end{proposition}
\begin{proof}
Put $p=f^{1/n}$, use the probability measure $\nu=A_0\wedge T=b^{1/n}p k_T\,\mathrm{d}V$ from the proof of \cref{prop:semipositive-dual-exponential-limit}, where the equality is \eqref{eq:semipositive-dual-trace-equality}, and set
\[
 q_\lambda=(c_\lambda/b)^{1/n}\mathrm{e}^{\lambda v_\lambda/n},
 \qquad R_\lambda=A_\lambda\wedge T-q_\lambda\nu\geq0.
\]
Here $R_\lambda\geq0$ by the trace inequality used in that proof, which also gives $c_\lambda\geq b$ and $\int_X(-v_\lambda)\,\mathrm{d}\nu\leq C_0$ for a constant $C_0$. Write $s_+=\max(s,0)$ for the positive part of a real number $s$. Since $q_\lambda\geq\mathrm{e}^{\lambda v_\lambda/n}$ and $1-\mathrm{e}^{x}\leq-x$, we obtain
\[
 (1-q_\lambda)_+\leq1-\mathrm{e}^{\lambda v_\lambda/n},
 \qquad
 \int_X(1-q_\lambda)_+\,\mathrm{d}\nu\leq C_0\lambda/n.
\]
Since $A_\lambda\wedge T$ and $\nu$ both have mass one,
\[
 \|A_\lambda\wedge T-\nu\|_{\mathrm{TV}}
 \leq R_\lambda(X)+\int_X|q_\lambda-1|\,\mathrm{d}\nu
 =2\int_X(1-q_\lambda)_+\,\mathrm{d}\nu.
\]
This proves the first estimate, with $C=2C_0/n$, and also gives
$R_\lambda(X)\leq C_0\lambda/n$ and
$\int_X|q_\lambda-1|\,\mathrm{d}\nu\leq2C_0\lambda/n$.

On $U=X\setminus Z$, write $T=\gamma^{n-1}$ and $A_0=a\gamma$ as in the proof of \cref{thm:semipositive-dual-exact-equation}, and put $\alpha_\lambda=q_\lambda a$. Decompose the positive current $A_\lambda$ into its absolutely continuous part, with matrix density $B_\lambda$ relative to $\gamma$, and its positive singular part. Since $\mathrm{MA}^{\mathrm{ext}}_\omega(v_\lambda)=c_\lambda\mathrm{e}^{\lambda v_\lambda}f\,\mathrm{d}V$, the trace inequalities for all Hermitian forms imply, as in the proof of \cref{cor:smooth-dual-exact-equation},
\[
 (\det B_\lambda)^{1/n}\geq\alpha_\lambda
 \quad\text{almost everywhere on }U.
\]
\Needspace{9\baselineskip}
This is a statement about determinants of matrices; no nonlinear product of coefficient measures is involved. Put $s_\lambda=\operatorname{tr}B_\lambda/n$ and $\delta_\lambda=s_\lambda-\alpha_\lambda\geq0$. On $U$, the absolutely continuous part of the trace error $R_\lambda$ is $\delta_\lambda\gamma^n$, and its singular part is the trace measure of the singular part of $A_\lambda$. Thus
\[
 \int_U s_\lambda\gamma^n\leq1,\qquad
 \int_U\delta_\lambda\gamma^n\leq C_0\lambda/n,
\]
and the singular trace mass on $U$ has the same upper bound.

For nonnegative numbers $\ell_1,\ldots,\ell_n$ with geometric mean at least $\alpha\geq0$ and arithmetic mean $s$, the arithmetic--geometric mean inequality and the Cauchy--Schwarz inequality give
\[
 \sum_i(\sqrt{\ell_i}-\sqrt\alpha)^2\leq n(s-\alpha),
 \qquad
 \sum_i|\ell_i-\alpha|\leq2n\sqrt{s(s-\alpha)}.
\]
Indeed $\sum_i\sqrt{\ell_i}\geq n\sqrt\alpha$, while
$\sum_i(\sqrt{\ell_i}+\sqrt\alpha)^2\leq4ns$.
Applying this to the eigenvalues of $B_\lambda$ relative to $\gamma$ and integrating, with the Cauchy--Schwarz inequality once more, we see that the absolutely continuous part of $A_\lambda-\alpha_\lambda\gamma$ has integrated trace norm at most $2n\sqrt{C_0\lambda/n}$ on $U$. Its singular part has norm at most $C_0\lambda$ in the same metric. Finally
\[
 \int_U|\alpha_\lambda-a|\gamma^n
   =\int_U|q_\lambda-1|\,\mathrm{d}\nu
   \leq2C_0\lambda/n.
\]
Since $\gamma$ is comparable to the fixed metric on the compact set $K\subset U$, these estimates and $\lambda\leq\sqrt\lambda$ for $0<\lambda\leq1$ give the second conclusion.
\end{proof}

The coefficient estimate is local on $X\setminus Z$. Convergence in capacity across $Z$ requires the additional envelope condition of \cref{cor:semipositive-dual-capacity-limit}.
